\documentclass[../main]{subfiles}
\begin{document}

\section{Mantenimiento Predictivo (parte 2) ultrasonido y análisis de lubricantes}

\subsection{Análisis con ultrasonido}
\subsubsection{Ultrasonido pasivo:}
Es producido por mecanismos rotantes, fugas de fluido, pérdidas de vacío, y arcos eléctricos. Pudiéndose detectarlo mediante la tecnología apropiada.
El Ultrasonido permite:
\begin{itemize}
\item Detección de fricción en maquinas rotativas.
\item Detección de fallas y/o fugas en válvulas.
\item Detección de fugas de fluidos.
\item Pérdidas de vacío.
\item Detección de "arco eléctrico".
\item Verificación de la integridad de juntas de recintos estancos.
\item Erosión.
\item Corrosión.
\item Pérdida de material cerámico en álabes o en placas aislantes.
\item Roces entre álabes fijos y móviles.
\item Decoloraciones en álabes del compresor, por alta temperatura.
\item Pérdidas de material de los álabes del compresor que se depositan en los álabes de turbina o en la cámara.
\item Deformaciones.
\item Piezas sueltas o mal fijadas, sobre todo de material aislante.
\item Fracturas y agrietamiento en álabes, sobre todo en la parte inferior que los fija al rotor.
\item Marcas de sobretemperatura en álabes.
\item Obstrucción de orificios de refrigeración.
\item Daños por impactos provocados por objetos extraños (FOD, Foreign object damages).

\end{itemize}
\subsection{Análisis de lubricantes}
\end{document}
